\documentclass[11pt]{article}
\usepackage[T1]{fontenc}
\usepackage[margin=1in]{geometry}
\usepackage{amsmath,amssymb,amsthm,mathtools,hyperref}
\hypersetup{hidelinks}
\newtheorem{theorem}{Theorem}[section]
\newtheorem{lemma}[theorem]{Lemma}
\newtheorem{proposition}[theorem]{Proposition}
\newtheorem{corollary}[theorem]{Corollary}
\theoremstyle{definition}
\newtheorem{remark}[theorem]{Remark}
\newcommand{\Q}{\mathbf Q}
\newcommand{\C}{\mathbf C}
\newcommand{\F}{\mathbf F}
\newcommand{\Z}{\mathbf Z}
\newcommand{\SL}{\operatorname{SL}}
\newcommand{\GL}{\operatorname{GL}}
\newcommand{\PGL}{\operatorname{PGL}}
\newcommand{\Sp}{\operatorname{Sp}}
\newcommand{\CH}{\operatorname{CH}}
\newcommand{\im}{\operatorname{im}}
\newcommand{\id}{\operatorname{id}}
\newcommand{\Alt}{\operatorname{Alt}}
\newcommand{\Hdg}{\operatorname{Hdg}}
\newcommand{\Hom}{\operatorname{Hom}}
\newcommand{\End}{\operatorname{End}}
\newcommand{\Gal}{\operatorname{Gal}}
\title{All powers of very general Jacobians of elementary abelian covers}
\author{Kyle Kistner}
\date{5 September 2026}
\begin{document}
\maketitle
\begin{abstract}
For every prime $p$, every elementary abelian group $(\Z/p)^m$, and every
 admissible branch datum, we give a proof of the rational Hodge conjecture
 for all powers of the very general Jacobian in the full branch family of
 covers of $\mathbf P^1$. For odd primes, the proof combines
 an explicit classification of connected monodromy blocks with a classical
 reduction of determinant classes to Fermat varieties and Shioda's theorem.
 A rational
 algebraic correspondence from determinant-word Hodge structures spans the
 monodromy invariants. Semisimplicity then lifts Hodge classes to balanced
 sources, where Fermat transfer applies. This avoids the incorrect identification of
 the full monodromy commutant with Hodge endomorphisms and requires no formula
 for the central Hodge torus. The binary case uses independent symplectic
 factors and divisor classes. A companion extends the argument to arbitrary
 finite abelian groups of odd prime-power exponent and certain mixed-prime
 exponents. This is a written proof draft using precisely
 stated published inputs and the companion finite-quotient proof; it is not a complete
 Lean formalization or an independently refereed result. Novelty is not asserted.
\end{abstract}

\section{The theorem and its scope}\label{sec:scope}
Let $p$ be prime, $G=(\Z/p)^m$ with $m\geq1$, and let
$a_1,\ldots,a_r\in G\setminus\{0\}$ generate $G$ and satisfy
$\sum_i a_i=0$. Assume $r\geq3$. Let $B=M_{0,r}$ be the full space of
ordered distinct branch points modulo projective transformations. The
specified local monodromies define a family of connected $G$-covers of
$\mathbf P^1$, interpreted on a finite cover when necessary to choose a
family and auxiliary graphs. Write $X_b$ for a fiber and $A_b=J(X_b)$.
The full branch family is essential: an arbitrary subfamily is not substituted
for $B$.

\begin{theorem}\label{thm:main}
There is a countable union $Z\subset B$ of proper closed algebraic subsets
such that, for every $b\in B(\C)\setminus Z$, every $q\geq1$, and every
$d\geq0$, the cycle class map
\[
 \CH^d(A_b^q)_\Q\longrightarrow
 H^{2d}(A_b^q,\Q)\cap H^{d,d}(A_b^q)
\]
is surjective. If $p=2$, these Hodge classes are generated by divisor classes.
\end{theorem}

For $r=2$ the only connected data generate a cyclic group and the curve has
genus zero, so the conclusion is immediate. For $r=3$ the parameter space is
a point; the odd-prime proof below includes its constant rank-one factors.
The theorem concerns rational classes over $\C$. It makes no claim about
integral classes, special fibers with additional Hodge classes, nonabelian
groups, or explicit equations for all cycles. Theorem~\ref{thm:main} itself
has prime exponent. The separate companion \cite[Section 7]{Extensions}
extends the conclusion to finite abelian groups of exponent dividing
$p^e$ for odd $p$, $3^a5^b$ or $3^a7^b$, and gives a criterion for any odd
exponent conditional on the corresponding Fermat-Jacobian Hodge conjecture.

The companion manuscript \cite{Fermat}, Theorem~1.1 and Corollary~1.2,
realizes each cyclic-cover determinant as an algebraic direct summand of
prime-degree Fermat cohomology. Shioda's product theorem \cite{Shioda}
then makes every rational Hodge class in arbitrary products of these
determinants algebraic. This includes all balanced mixed determinants and
requires no monodromy hypothesis. The construction has direct precedent
in Schoen's 1989 work, as recorded in \cite{Fermat}; its algebraicity
consequence is not presented as a new cycle-construction theorem.
The earlier companion \cite{Fusion} supplies an alternative degeneration
argument, but is no longer a necessary premise. The original frozen Phase I
and Phase II files are not premises either.

There is also substantial monodromy precedent. Rohde
\cite[Construction 3.2.1, Theorem 5.3.7]{Rohde} treats arbitrary branch
residues for prime cyclic covers and proves their full derived monodromy.
Together with the classical determinant construction and invariant theory,
this already implies the cyclic case of Theorem~\ref{thm:main}; that
implication is our comparison, not a claim that Rohde states the Hodge
conclusion verbatim. Spelta--Tamborini \cite[Theorem 1.7]{ST} treat abelian
covers with pairwise nonisomorphic noncompact unitary factors, and explicitly
study the diagonal obstruction when factors repeat. The possible additional
scope here is the uniform classification for distinct character words,
including all repeated-factor cases and their actual Hodge orientations.
The methods of comparing Dehn-twist spectra and subdirect products have
these precedents. The four-point Klein-four cover symmetries themselves
are classical; see also \cite[Section 2.6, Lemma 16]{FMZ} in the square-tiled
description. Section~\ref{sec:monodromyexamples} exhibits the additional
scope and the genuine graph exceptions concretely. No priority for the
assembled theorem is asserted.

\section{Character factors and a mixed embedding}\label{sec:moving}
Suppose first that $p$ is odd. Put $K=\Q(\zeta_p)$ and
$\Gamma=\F_p^\times=\Gal(K/\Q)$. The evaluation code is
\[
 C=\{(c(a_1),\ldots,c(a_r)):c\in\Hom(G,\F_p)\}\subset\F_p^r.
\]
Generation of $G$ makes evaluation injective. We use the same symbol $c$
for a character and its word. Write $S(c)=\{i:c_i\ne0\}$,
$s(c)=|S(c)|$, and $h(c)=s(c)-2$. Over $K$ or $\C$, let $V_c$ be the
$\zeta_p^{c(g)}$ character part of $V=H^1(A_b,\Q)$. A uniform reversal
of this convention has no effect on the conclusions.

The quotient $Y_c=X_b/\ker c$ is a cyclic $p$-cover with active residues
$(c_i)_{i\in S(c)}$. Its positive-genus rational factor
$U_c=H^1(JY_c,\Q)$ has $K$-rank $h(c)$, and
\[
 U_c\otimes_\Q\C=\bigoplus_{t\in\Gamma}V_{tc}.
\]
One representative from each nonzero character line gives the rational
isogeny decomposition of $A_b$. Pullback and norm realize the associated
rational cohomological inclusions and projections algebraically. Support two
has genus zero; support one is impossible by the sum-zero condition. We use
only supports at least three below.

\begin{lemma}[A mixed embedding always exists]\label{lem:mixed}
For nonzero residues $A_1,\ldots,A_s\in\F_p$ with sum zero and $s\geq4$,
there is $t\in\Gamma$ such that
\[
 2\leq Q(t):=\frac1p\sum_i[tA_i]_p\leq s-2.
\]
Consequently both Hodge multiplicities of the corresponding character
space are positive at some embedding.
\end{lemma}
\begin{proof}
The integers $Q(t)$ lie between $1$ and $s-1$, and
$Q(p-t)=s-Q(t)$. If $Q(1)$ is already interior there is nothing to prove.
Otherwise negate the word to arrange $Q(1)=1$, and use positive representatives
$A_i$ with $\sum A_i=p$. For $1\leq t\leq p-2$,
\[
 Q(t+1)-Q(t)=1-\#\{i:[tA_i]_p+A_i\geq p\}\leq1.
\]
No endpoint residue is zero in this range. Since $Q(p-1)=s-1\geq3$,
the first crossing of $2$ equals $2$. The two Hodge dimensions are
$Q(t)-1$ and $s-1-Q(t)$, in one order, by the cyclic eigendifferential
formula \cite[Theorem~2.1]{MN}; see also \cite[Section~6]{Fusion}.
\end{proof}

\begin{proposition}\label{prop:projections}
Let $M$ be the connected algebraic monodromy group on $V$. Every complex
projection of $M$ on $V_c$ with $s(c)\geq4$ is $\SL_{h(c)}$. This includes
the definite conjugate embeddings. The only nonzero character factors with
trivial connected monodromy have support three and rank one.
\end{proposition}
\begin{proof}
Choose an active point as infinity and set $n=s-1$. At the embedding of
Lemma~\ref{lem:mixed}, the sum of finite fractional weights is
$Q(t)-[tA_\infty]_p/p$, strictly between $1$ and $n-1$. All finite
weights are active, $n\geq3$, and the gcd condition is automatic at a prime.
Thus Definition~1.1(a) and Theorem~5.1 of \cite{MN} give full special
unitary projection, whose complexification is $\SL_{s-2}$, on the affine
configuration space $\operatorname{Conf}_{s-1}(\C)$. Normalizing two finite
points to $0,1$ identifies that space with
$\operatorname{Aff}_1\times M_{0,s}$. The change $x=\alpha z+\beta$
rescales the Kummer coordinate by
$\alpha^{(\sum_{i\ne\infty}[A_i]_p)/p}$; thus the affine-group fiber
has only finite scalar monodromy on each character. Passing to $M_{0,s}$
preserves the connected closure. Finally $M_{0,r}\to M_{0,s}$ has
connected configuration-space fibers and surjects on fundamental groups.
This transfers the conclusion to the full branch family without asserting
a surjection from $\pi_1(M_{0,s})$ to an affine pure braid group.

In a $K$-basis of the eigenspace in $V\otimes_\Q K$, monodromy matrices
have entries in $K$. The Zariski closure of a set of $K$-points commutes with
extension from $K$ to $\C$: expand the coefficients of a vanishing complex
polynomial in a finite $K$-linearly independent list and evaluate at the
$K$-points. Each resulting $K$-coefficient polynomial vanishes separately.
Applying $\zeta_p\mapsto\zeta_p^u$ therefore transports the closure to
that on $V_{uc}$. Full $\SL_h$ at one embedding gives full $\SL_h$ at
every embedding. Pair-braid determinants and the deck scalar ambiguities
have finite order \cite[Corollary~2.7]{MN}, so there is no connected scalar
factor in this argument.

With three active points, normalization to $0,1,\infty$ makes the cyclic
family isotrivial after a finite root choice. Its monodromy is finite.
The positive-rank cases have now all been accounted for. A definite complex
eigenspace alone has no rational integral lattice; compactness at that one
embedding does not contradict its full complex Zariski closure.
\end{proof}

\section{Exact monodromy blocks and actual oriented graphs}\label{sec:blocks}
For four labelled positions $S$, let $\mathcal V_4(S)$ be the Klein four
group of double transpositions and the identity. Orbit notation below always
means intersection with the evaluation code $C$; $C$ need not be preserved
by all these permutations.

\begin{theorem}\label{thm:blocks}
On the indices $c\in C$ with $s(c)\geq4$, connected monodromy has the
following complex simple blocks:
\begin{enumerate}
\item Different active supports belong to independent factors.
\item On a fixed support of size at least five, a block is $\{c,-c\}$,
with standard and dual representations of $\SL_{h(c)}$.
\item On a support of size four, a block is
$\bigl(\mathcal V_4(S)c\cup-\mathcal V_4(S)c\bigr)\cap C$.
Its oriented component $O_c=(\mathcal V_4(S)c)\cap C$ consists of
actual algebraically isomorphic character spaces; $-O_c$ is its polarized
dual component. The two components may coincide.
\end{enumerate}
The connected linear image is the product of these $\SL_h$ factors with
the indicated occurrences. The entire description is Galois-stable.
\end{theorem}

\subsection{The same projective group before finite base change}
Projective character monodromy is defined on the same group
$\Pi=\pi_1(M_{0,r})$ for every $c$: two lifts of a marked-sphere mapping
class differ by a deck transformation, hence by a scalar on $V_c$.
This gives a product of projective representations. It does not claim one
common scalar on their direct sum. Root choices or level covers restrict to
finite-index subgroups and leave connected closures unchanged.

For a centerless simple algebraic group $H$ and an automorphism $\theta$,
\begin{equation}\label{eq:normalizer}
 N_{H\times H}(\operatorname{Graph}(\theta))
 =\operatorname{Graph}(\theta).
\end{equation}
Indeed a normalizing pair $(a,b)$ makes $\theta(a)^{-1}b$ centralize $H$.
The connected component is normal in the full closure. Consequently, when a
pair of connected projective projections is a graph, \emph{every original
braid image} belongs to that graph. We may compare finite-order pair braids
before finite base changes raise their eigenvalues to powers.

The connected solvable radical of $M_\C$ projects to a connected solvable
normal subgroup of each full $\SL_h$, hence trivially. It also acts
trivially on the finite-monodromy summands. Faithfulness makes the radical
trivial. Semisimple Goursat theory now partitions the projective projections
into graph blocks. At the Lie-algebra level, each simple target receives
exactly one simple ideal: two commuting nonzero ideals cannot both surject
onto it. In particular there are no extra relations invisible in pairwise
projections.

\subsection{Nonidentity twists, including zero pair sums}
Let $i,j$ be active for $c$, with $s(c)\geq4$, and choose another active
label as infinity. Put $n=s(c)-1$; the finite residue sum is nonzero and
$\dim V_c=n-1\geq2$. Number the coordinate description so the chosen
pair is $(1,2)$; this relabels a description, not an action by a pure braid.
Theorem~2.2 of \cite{MN} gives $n$ spanning vectors $g_1,\ldots,g_n$
with their only relation $\sum_i g_i=0$. Thus $g_1\ne0$. Theorem~2.5
there gives
\[
 \rho_c(A_{12})-I=\kappa\,g_1\otimes\langle -,g_1\rangle,
 \qquad \kappa\ne0,
\]
where the Hermitian form is nondegenerate and the coefficient is nonzero
because both local residues are active. Hence the difference has rank one
and is nonzero. It is not a scalar difference in dimension at least two.
Every active pair twist is therefore projectively nontrivial.

For the chosen positive-eigenvalue character convention, put $q=\zeta_p$
in the source pair formula. If $c_i+c_j\ne0$, its spectrum is
\begin{equation}\label{eq:reflection}
 (\zeta_p^{c_i+c_j},1,\ldots,1).
\end{equation}
If the sum is zero, the same rank-one argument proves it is a
\emph{nonidentity} unipotent. This last conclusion cannot be inferred from
the word ``unipotent'' alone. For example, at $p=3$ the zero-sum word
$(1,2,2,2,2)$ has zero sum on every pair containing its first position.

If $S(c)\ne S(e)$ and their ranks agree, choose $i\in S(c)\setminus S(e)$
and $j\in S(c)\setminus\{i\}$. The same pair twist is nontrivial on
$V_c$, but becomes trivial projectively on $V_e$ after forgetting $i$.
This contradicts \eqref{eq:normalizer}. Unequal ranks already give
nonisomorphic simple groups. This proves support separation.

\subsection{Reconstruction on a fixed support}
For $h\geq3$, every complex automorphism of $\PGL_h$ is inner or inner
followed by inverse transpose. A graph therefore has a single fixed sign
$\epsilon\in\{1,-1\}$. The multiplicities $1,h-1$ in
\eqref{eq:reflection} distinguish the exceptional eigenvalue projectively.
Comparison of all finite pair braids, including the unipotent case, gives
\[
 e_i+e_j=\epsilon(c_i+c_j).
\]
For three distinct finite indices,
\[
 2e_i=(e_i+e_j)+(e_i+e_k)-(e_j+e_k)=2\epsilon c_i.
\]
Since $p$ is odd, all finite coordinates agree up to $\epsilon$, and the
sum-zero relation gives the infinity coordinate. Thus $e=\epsilon c$.
The converse follows from identity and polarization.

For $h=2$, write the four active coordinates as $c_1,c_2,c_3,c_4$ and put
$x=c_1+c_2$, $y=c_1+c_3$, $z=c_1+c_4$. The inverse is
\begin{equation}\label{eq:signcube}
 c=\tfrac12(x+y+z,\ x-y-z,\ -x+y-z,\ -x-y+z).
\end{equation}
All automorphisms of $\PGL_2$ are inner. The unordered eigenvalue ratios
of the three finite pair braids give independent signs on $x,y,-z$.
Thus $(x_e,y_e,z_e)=(\epsilon_x x,\epsilon_y y,\epsilon_z z)$.
The four even sign changes are exactly the effects of
$\id,(12)(34),(13)(24),(14)(23)$; the other four are their negatives.
Equation~\eqref{eq:signcube} proves $e=\epsilon Pc$. Zero pair sums only
enlarge stabilizers and give no additional solutions.

For the geometric converse, normalize the four branch points to
$0,1,t,\infty$. The corresponding double transpositions are
\[
 T_{12|34}(x)=\frac{t(x-1)}{x-t},\quad
 T_{13|24}(x)=\frac{x-t}{x-1},\quad
 T_{14|23}(x)=\frac{t}{x}.
\]
Write the cyclic equations as $y^p=f_c(x)$ and $z^p=f_e(x)$.
For $e=Pc$, the divisor of $f_e(T_P(x))/f_c(x)$ is divisible by $p$.
A degree-zero divisor on $\mathbf P^1$ is principal, so
\[
 f_e(T_P(x))/f_c(x)=a(t)h(x)^p.
\]
Here $a(t)$ is a nonzero constant times powers of $t$ and $t-1$; this
also follows directly from the displayed transformations. After adjoining
$a(t)^{1/p}$, a finite etale cover of $M_{0,4}$, the map
\[
 (x,y)\longmapsto(T_P(x),a(t)^{1/p}h(x)y)
\]
is deck-equivariant and extends to an isomorphism of the smooth projective
curves. Its graph, or inverse graph as appropriate for cohomological
variance, identifies $V_{bc}$ with $V_{be}$ for \emph{every} $b\in\Gamma$.
For the full family pull it back along the active-point forgetting map.
Pullback and norm for the cyclic quotients transfer these graphs to the
rational factors of $A_b$.

The negative component is obtained using the polarized identification
$V_{-c}\simeq V_c^*(-1)$, together with these oriented graphs. At the
level of $\SL_2$ representations the dual is isomorphic to the standard,
which gives the claimed adjoint graph. It does not assert a degree-zero
Hodge isomorphism between an arbitrary $V_c$ and $V_{-c}$.
If $O_c=-O_c$, a double transposition sends $c$ to $-c$; its two pairs
are opposite residues, and every embedding has signature $(1,1)$. The
oriented graph to the negative character is then actual, as asserted.

These arguments give all necessary and sufficient pairwise graph relations.
Goursat supplies the product. Finite discrepancies between linear lifts
vanish on the connected group; its Lie algebra and the standard/dual
representations lift to the product of the simply connected groups
$\SL_h$. Each factor has a faithful standard occurrence, so its connected
linear image is exactly $M_\C$. Multiplying all residues by a unit
preserves every displayed relation. This completes
Theorem~\ref{thm:blocks}.

\begin{remark}[A false endomorphism shortcut]\label{rem:false}
At $p=5$, the word $(1,1,1,2)$ has signatures $(2,0),(1,1),(1,1),(0,2)$,
up to convention. Its opposite character is isomorphic as an $\SL_2$
module, but at a definite embedding has the opposite pure Hodge type.
The nonzero monodromy Hom there is not a Hodge Hom. We do not identify
$\End_M(V)$ with the Hodge endomorphism algebra in general.
\end{remark}

\section{Determinant words and algebraic rational projectors}\label{sec:words}
For each positive-genus character $c$, choose a deck generator on $Y_c$
whose induced action has eigenvalue $\zeta_p^b$ on $V_{bc}$. Equivalently
choose $g_c\in G$ with $c(g_c)=1$. Different scalar multiples of $c$
may use the same curve with different generators; they remain distinct
word entries. A nonempty effective word $w=(c_1,\ldots,c_N)$ permits
repetitions. Put $H=\sum_j h(c_j)$. Its rational determinant source is
\begin{equation}\label{eq:word}
 W_w\otimes_\Q\C=
 \bigoplus_{b\in\Gamma}\bigotimes_{j=1}^N
                \bigwedge^{h(c_j)}V_{bc_j}
 \ \subset\ H^H\left(\prod_j JY_{c_j},\C\right).
\end{equation}
It has rational dimension $p-1$ and rank one over $K$. This is the
matching-embedding tensor product over $K$, not the full tensor product
of rational determinant spaces.

We recall explicitly why its realization and the operations used below
are rational algebraic. On the full Jacobian the cyclic-line projector is
\[
 \Pi_c=\frac1{|\ker c|}\sum_{g\in\ker c}g^*
       -\frac1{|G|}\sum_{g\in G}g^*.
\]
Its image is $U_c$; the subtracted invariant $H^1$ is zero here. On this
image put $T_c=g_c^*$. For raw $H^1$ slots with intended character
labels $c_1,\ldots,c_n$, first apply their $\Pi_{c_i}$ and then
\begin{equation}\label{eq:matching}
 E(c_1,\ldots,c_n)=
 \prod_{i=2}^n\left(\frac1p\sum_{a=0}^{p-1}
           T_{c_1,1}^aT_{c_i,i}^{-a}\right).
\end{equation}
A root-of-unity sum shows that its image is exactly
$\bigoplus_b\bigotimes_i V_{bc_i}$. Each term is a rational algebraic
deck correspondence. On identical slots the matching projector commutes
with ordinary permutations, since equality of all embedding labels is
permutation-invariant.

The degree projectors on an abelian variety are cohomologically algebraic:
for any integer $m\geq2$, interpolate the distinct eigenvalues $m^i$ of
$[m]^*$ to obtain
\[
 \pi_k=\prod_{\substack{0\leq i\leq2\dim D\\i\ne k}}
               \frac{[m]^*-m^i}{m^k-m^i}.
\]
For addition $\mu:D^n\to D$, diagonal $\delta:D\to D^n$, and
$P=\pi_1^{\otimes n}$, set $\iota_n=P\mu^*$ and $q_n=\delta^*P$.
On the relevant degree summands,
\[
 q_n\iota_n=n!\id,\qquad \iota_nq_n=n!\Alt_n.
\]
Here $\Alt_n$ is the ordinary antisymmetrizer. Ordinary permutations
are realized by geometric factor permutations with their degree-one Koszul
sign corrected. Thus insertion of alternating volume tensors, and passage
back to exterior cohomology, are algebraic correspondences with explicit
nonzero rational normalizations.

Apply \eqref{eq:matching} to all raw word slots, alternate within each
volume block, and return using $q_{h(c_j)}$ and the factorials. This
constructs an algebraic rational projector with image \eqref{eq:word}.
On the matched image, action by $\zeta_p$ on \emph{one raw slot} acts
by $\zeta_p^b$ on its $b$-line and commutes with the alternating
restrictions. It gives a faithful $K$-action by algebraic Hodge
endomorphisms on $W_w$. Simultaneous deck action on all $h(c_j)$ slots
would instead give $\zeta_p^{b h(c_j)}$, possibly $1$ when $p\mid h(c_j)$;
that incorrect action is not used. These constructions are also detailed
in \cite[Section~2]{Fusion}.

\begin{proposition}[The source Hodge criterion]\label{prop:balance}
Suppose $H+2t=2d$, where $t\geq0$. Then the rational Hodge-class space
of $W_w\otimes\Q(-t)$ is either zero or the whole source. It is the
whole source precisely when
\begin{equation}\label{eq:balance}
 2\sum_{j=1}^N\sum_{i\in S(c_j)}[b c_{j,i}]_p
       =p\sum_{j=1}^N s(c_j)\qquad(b\in\Gamma).
\end{equation}
In that case every source Hodge class is algebraic, in its specified
cohomological realization with the Tate factor.
\end{proposition}
\begin{proof}
The rational Hodge-class subspace is stable under the $K$-action by Hodge
endomorphisms. Since $W_w$ is a $K$-line, a nonzero such subspace is the
whole line. Equivalently, a nonzero rational vector has nonzero components
at all embeddings of $K$.

Let $d_c(b)=\dim V_{bc}^{1,0}$. The $b$-line in \eqref{eq:word} has
type $(\sum_jd_{c_j}(b),H-\sum_jd_{c_j}(b))$. After twisting by
$\Q(-t)$ it is of type $(d,d)$ exactly when
$\sum_jd_{c_j}(b)=H/2$. The cyclic dimension formula gives
\eqref{eq:balance}; the opposite convention merely replaces $b$ by $-b$.
Here $H$ is even since $H+2t=2d$. The Fermat transfer
\cite[Corollary~1.2]{Fermat} applies to this exact word, including repeated
curves, rank-one factors and chosen generators. Together with Shioda's
product theorem it makes all of $W_w$ algebraic. The standalone fusion
proof \cite[Theorem~1.1]{Fusion} is an alternative route.
Tate factors are represented by products of point or divisor
classes. No genericity is required for this last application.
\end{proof}

\section{Rational diagrams span every monodromy invariant}\label{sec:diagrams}
We prove the generator step rather than treating it as a hypothesis about
the central torus. All arguments are in a fixed finite tensor degree.

\begin{lemma}[Mixed tensor invariants]\label{lem:fft}
Let $U$ have dimension $h\geq2$ over a characteristic-zero field.
Invariants for $\SL(U)$ in tensor products of standard and dual copies
are spanned by diagrams made from coevaluations, alternating $h$-slot
standard volumes, and alternating $h$-slot dual volumes.
\end{lemma}
\begin{proof}
We may extend scalars to an algebraic closure, since invariants and the
claimed linear span commute with extension. For $a$ standard and $b$ dual
slots, the center $\mu_h$ forces
$a-b$ to be divisible by $h$. If $a=b+mh$ with $m\geq0$, the invariant
space, after choosing a nonzero volume in $\det(U)^m$, identifies with
\[
 \Hom_{\SL(U)}(U^{\otimes b},U^{\otimes a})
 \simeq\Hom_{\GL(U)}(U^{\otimes b}\otimes\det(U)^m,U^{\otimes a}).
\]
The scalar degrees on the latter source and target agree, so $\SL$
equivariance gives $\GL$ equivariance. Alternating $m$ disjoint blocks
of $h$ slots embeds that source as a $\GL$-equivariant direct summand
of $U^{\otimes a}$. Extend a map using the alternating projector.
Schur--Weyl duality \cite[Theorem~4.57(i), Proposition~4.58]{Etingof} says
every $\GL$-equivariant endomorphism of this tensor power is a linear
combination of slot permutations. This produces
$m$ volumes and $b$ coevaluations. If $a<b$, interchange $U$ and $U^*$.
The argument is valid also when $h<a$; permutation operators need span,
not be linearly independent.
\end{proof}

Fix $q\geq1$ and $n=2d$. Apply the lemma to the product of blocks in
Theorem~\ref{thm:blocks}, with occurrences from $H^1(A_b^q)$ labelled
separately. Within an oriented component, use its actual graph maps to
identify the standard copies. The negative component is the polarized dual
$U^*(-1)$, so a standard--dual coevaluation has source $\Q(-1)$ and
is a divisor tensor. A same-orientation volume has source $\det U$;
it is not assumed Hodge. Dual volumes use the actual negative-character
determinant, of weight $h$, including its Tate bookkeeping.

For a self-dual rank-two component $O_c=-O_c$, one may identify every
copy as standard by the actual graphs and use two-slot volumes. Keeping
them as determinant sources is harmless. Support-three factors have rank
one and trivial connected monodromy; each of their slots is itself a
one-slot determinant source. Thus no constant CM part is omitted.
Every diagram has a nonempty word $w$ of total volume degree $H$ and
$t$ pairing blocks with $H+2t=n$, or has an empty word and consists
entirely of pairings. All multiplicities are nonnegative.

\begin{lemma}[Algebraic rational diagram realization]\label{lem:rational}
For each such complex diagram and its full Galois orbit there is a rational
algebraic Hodge map into $H^n(A_b^q,\Q)$ from $W_w\otimes\Q(-t)$,
or from a finite sum of $\Q(-t)$ for an empty word, whose scalar
extension contains every diagram in that orbit in its image. Collectively
these maps have image exactly $H^n(A_b^q,\Q)^M$.
\end{lemma}
\begin{proof}
First work in raw tensor slots, with intended labels $c_1,\ldots,c_n$.
For each volume choose its representative character $a_j$ and its cyclic
quotient $Y_{a_j}$. The alternating inclusion from its determinant source,
followed by the actual oriented graphs and quotient pullbacks, maps its
$b$-line to the prescribed slots $V_{bc_i}$. The maps work simultaneously
for every $b$; no individual complex eigenprojector is called rational.

For each pairing use the rational polarization coevaluation on the
corresponding cyclic factor. Its extension is the sum of the nonzero
coevaluations in $V_{ba}\otimes V_{-ba}$, for $b\in\Gamma$. A
Poincare divisor, algebraic degree projections, and rational inverse
isogeny for the polarization realize this tensor algebraically. Apply the
actual oriented graphs to the two slots. This produces precisely the
conjugates of the desired standard--dual pairing with nonzero compatible
normalizations.

For a nonempty word, append the $t$ rational pairings to $W_w$, reorder
the slots, and apply the \emph{global} projector
$E(c_1,\ldots,c_n)$ from \eqref{eq:matching}. Before this last operation
the pairing blocks may have independent embedding labels. The projector
forces all their labels to equal the word label. Consequently we obtain a
rational algebraic Hodge map
\[
 F_{\mathcal D}:W_w\otimes\Q(-t)
       \longrightarrow H^1(A_b^q,\Q)^{\otimes n}
\]
whose restriction to the $b$-line is the $b$-conjugate diagram. It is
nonzero on that tensor diagram before the exterior projection. Scalar
extension separates the individual embedding lines, so its image contains
each of them.

For an empty word with $n>0$, apply the same matching projector to the
external product of the $t$ rational pairings. This gives an algebraic
tensor equal to the sum of the conjugate diagrams. Apply
$T_{c_1,1}^a$ for $0\leq a\leq p-2$. These are $p-1$ algebraic
maps from $\Q(-t)$. Their matrix on the $b$-diagrams is
$(\zeta_p^{ab})_{b\ne0,\,0\leq a\leq p-2}$, which is invertible:
a polynomial of degree at most $p-2$ vanishing at all $p-1$ primitive
$p$th roots is zero. Thus the full orbit, not just its trace, is spanned.
For $n=0$, use the unit $\Q\to\Q$.

All slot operations are algebraic by Section~\ref{sec:words}. Composing
with the exterior/cup-product projection gives the required target maps.
Diagrams that vanish or become dependent there do not matter: the
antisymmetrizer supplies an equivariant section of the exterior quotient,
so invariant tensor diagrams still span its invariant quotient.

There are only finitely many labels, partitions and slot permutations in
degree $n$. Include a finite spanning collection and its complete Galois
orbits. Each source is fixed by connected monodromy (determinants of
$\SL_h$, constant rank-one factors, and pairings), and the chosen graph
maps commute with monodromy after a common finite base change. The images
therefore lie in the invariants. Conversely the tensor lemma spans them
after extension to $\C$. Invariants and images commute with scalar
extension in these finite-dimensional rational representations, giving
exact equality over $\Q$.
\end{proof}

\begin{remark}
An abstract subdirect-product classification without the actual orientations
would not suffice. Take a common $\SL_2$ standard module $U$ with four
distinct central twists $\chi_i$ satisfying $\sum_i\chi_i=0$ but no
opposite pair. The product of alternating pairings in four different copies
can be invariant, while individual determinants and actual Hodge-Hom maps
have even signed multiplicity in each copy and cannot produce it. The
geometric graph construction in Section~\ref{sec:blocks} excludes these
unrelated twists within an oriented component. We never promote an
arbitrary projective monodromy graph to a Hodge correspondence.
\end{remark}

\section{Lifting Hodge classes and completing the odd-prime proof}\label{sec:lift}
The preceding construction yields a finite rational algebraic Hodge map
\begin{equation}\label{eq:surjection}
 F:\mathcal S=
 \left(\bigoplus_{\mathcal D}W_{w_{\mathcal D}}\otimes\Q(-t_{\mathcal D})\right)
 \oplus\left(\text{Tate summands}\right)
 \longrightarrow H^{2d}(A_b^q,\Q),
\end{equation}
with image $I=H^{2d}(A_b^q,\Q)^M$. Every summand has weight $2d$.
The image is a rational Hodge substructure because $F$ is a Hodge map;
this conclusion does not require a separate assertion about the Hodge-group
center.

\begin{lemma}[Hodge lifting]\label{lem:lift}
Every rational Hodge class in $I$ has a rational Hodge preimage in
$\mathcal S$.
\end{lemma}
\begin{proof}
Polarizable pure rational Hodge structures are semisimple. Concretely, a
polarization makes the orthogonal complement of $\ker F$ a Hodge
complement on which $F$ is an isomorphism onto $I$. Its inverse lifts the
class. This splitting is only a Hodge splitting. No algebraicity of the
splitting is assumed or used.
\end{proof}

At a very general point, all rational Hodge classes on all powers lie in
connected monodromy invariants. Here is the precise genericity input.
The standard generic monodromy inclusion is \cite[Section~4, Lemma~4]{Andre}.
For each power and degree, twist $R^{2d}f_*\Z$ by $(d)$ to obtain the
weight-zero variation (modulo torsion) with its polarization.
Algebraicity of the exceptional Hodge loci follows from
\cite[Theorem~1.1 and Corollary~1.2]{CDK}: bounded integral Hodge classes
form an algebraic space finite over the base. Remove the proper images of
its components, for every bound and each cohomological degree of each
power. There are countably many such proper closed algebraic subsets.
At the remaining points any Hodge class is generic and its monodromy orbit
is finite; equivalently, connected monodromy fixes it. One can see the
finiteness directly from the positive polarization on a locally constant
type-$(0,0)$ integral Hodge subsystem. Finite auxiliary base covers preserve
this conclusion and the meaning of very general after taking images of
the exceptional sets.

Let $\alpha\in H^{2d}(A_b^q,\Q)$ be a Hodge class at such a point.
It belongs to $I$. Lift it by Lemma~\ref{lem:lift} to a Hodge vector in
\eqref{eq:surjection}. Its components are Hodge vectors separately in
the displayed rational direct summands. By Proposition~\ref{prop:balance},
a nonzero component in a nonempty word forces that entire word to satisfy
all-row balance, and Fermat transfer supplies algebraic representatives. The Tate
components are algebraic too. All \emph{forward} maps in $F$ are
algebraic correspondences, external products with algebraic classes, and
pullbacks or pushforwards. They therefore carry those representatives to
algebraic cycles with class $\alpha$. This proves
Theorem~\ref{thm:main} for odd $p$.

Relations and cancellations among complex diagram monomials cause no gap:
we lift through a rational Hodge surjection, rather than assuming an
arbitrary complex expansion of $\alpha$ is termwise balanced. The proof
also uses all effective words needed in each fixed degree, not an integer
basis of a lattice that might require negative multiplicities.

\section{The binary case}\label{sec:binary}
For $p=2$, a nonzero character is determined by its support, whose size
$s$ is even. Its cyclic quotient is hyperelliptic of genus
$g=(s-2)/2$. Supports of size two vanish; the remaining ones have $s\geq4$.
Their rational $H^1$ factors have connected monodromy $\Sp_{2g}$.
For precise application of A'Campo \cite[Theorem~1, p.~319]{ACampo},
put all $s=2g+2$ active branch points in the affine line. His parameter
$\mu=s-1=2g+1$ is odd and at least three, including $g=1$.
The symplectic module is the homology of the affine curve modulo its
intersection radical, hence the homology of the compact curve.
The full braid image contains the principal level-two subgroup. Passing
to ordered configurations is finite index, so the connected Zariski closure
remains $\Sp_{2g}$. Passing from affine ordered configurations to $M_{0,s}$
has only finite deck-scalar fiber monodromy: lifts of the same marked-sphere
map differ by a deck transformation. Thus normalization also preserves the
connected closure. Pullback to $M_{0,r}$ uses the connected fibers of the
active-point forgetting map. No claim about equality of the discrete image after
arbitrary finite base change is needed.

Distinct supports give independent factors. Indeed the pair twist for two
active points has a nonzero Picard--Lefschetz transvection (squared for
ordered points). Its vanishing cycle is nonseparating: the double cover of
a disk containing those two branch points is an annulus; the double cover
of the complementary disk is connected because it contains the remaining
$s-2\geq2$ branch points. Cutting the annular core therefore leaves the
whole covering surface connected. Thus its homology class is nonzero.
If one of the two points is inactive for a different support, that twist
becomes projectively trivial on the latter quotient by forgetting it.
The normalizer argument \eqref{eq:normalizer}, now for the adjoint
symplectic group, and semisimple Goursat theory rule out a diagonal graph.
Different genera already give nonisomorphic simple groups.

The symplectic tensor invariant theorem
\cite[Section~7.3, Theorem~7.4 and Corollary~7.5]{Goodman} says invariants
in all tensor powers of standard modules are spanned by products of pairings.
For this
product of independent symplectic factors, each pairing is algebraic: use
the quotient polarization and its Poincare divisor, pulled back to the
chosen copies of $A_b$. The same algebraic exterior projection gives all
invariants in $H^*(A_b^q)$. Very-general Hodge classes are invariant by the
argument in Section~\ref{sec:lift}; hence they are generated by divisor
classes. This completes the binary case.

\section{Examples with repeated signature factors}\label{sec:monodromyexamples}
The block theorem determines both independence and actual graph relations
when signatures alone do not distinguish the factors. The following
examples concern different cyclic quotient lines, rather than different
embeddings of one cyclic quotient.

\subsection{Equal signatures and independent factors at $p=3$}
Take $G=\F_3^2$ and use the columns of
\[
 B=\begin{pmatrix}
 1&1&1&2&2&2\\
 1&1&2&1&2&2
 \end{pmatrix}
\]
as the six ordered branch vectors. They are nonzero, generate $G$, and
sum to zero. Let $c,e$ be the first and second rows. Both have support
six and $K$-rank four. For each of the two nonzero multipliers their
residue sum is $9$, so both signatures are $(2,2)$.

The rows are nonproportional. On the pair twist at labels $1,3$, their
residue sums are respectively $2$ and $0$. Its action on $V_c$ is
projectively semisimple of order three, whereas on $V_e$ it is a
nonidentity unipotent. The normalizer argument therefore excludes a
graph between these two equal-signature factors.

The other two projective character lines are represented by
\[
 c+e=(2,2,0,0,1,1),\qquad c-e=(0,0,2,1,0,0).
\]
They have support four and support two, respectively. Thus
Theorem~\ref{thm:blocks} gives
\[
                       M_\C=\SL_4\times\SL_4\times\SL_2.
\]
The full covering curve has genus ten, since
$2g-2=9(-2+6\cdot\frac23)=18$. Its Jacobian has cyclic quotient
dimensions $4,4,2$, with the fourth quotient of genus zero. The
nonrepetition hypothesis of \cite[Theorem~1.7]{ST} does not apply to the
two identical $\operatorname{SU}(2,2)$ signatures; their independence
here follows from the labelled pair operators.

\subsection{A genuine graph between different quotient lines at $p=5$}
Take $G=\F_5^2$ with branch matrix
\[
 B=\begin{pmatrix}1&1&1&2\\1&1&2&1\end{pmatrix},
 \qquad c=(1,1,1,2),\quad e=(1,1,2,1).
\]
Again the columns are nonzero, span $G$, and sum to zero. The rows are
nonproportional, but
\[
                         e=(12)(34)c.
\]
Consequently Theorem~\ref{thm:blocks} identifies their rank-two
character spaces by actual oriented graph correspondences. With the
branch points normalized to $0,1,t,\infty$, the map of the base is
\[
                         x\longmapsto\frac{t(x-1)}{x-t}.
\]
The Kummer lift after the finite root choice constructed in
Section~\ref{sec:blocks} identifies $V_{bc}$ with $V_{be}$ for every
$b\in\F_5^\times$. Thus different cyclic quotient lines can give a
genuine diagonal monodromy block. Equal signatures do not themselves
decide between this example and the preceding independent one.

\subsection{The genus-five example of Spelta--Tamborini}
We use exactly the ordered datum of \cite[Section~3, Example~2]{ST}.
Write its generators additively as $g_1,g_2,g_3$ and identify them with
the standard basis of $G=\F_2^3$. Their multiplicative notation for the
branch tuple then becomes
\[
 (g_2,g_2,g_2,g_1+g_2,g_1+g_2+g_3,g_2+g_3).
\]
In the same order of the six branch labels, the matrix is
\[
 B=\begin{pmatrix}
 0&0&0&1&1&0\\
 1&1&1&1&1&1\\
 0&0&0&0&1&1
 \end{pmatrix}.
\]
Let $c_i$ be its rows; the character corresponding to
$u=(u_1,u_2,u_3)$ evaluates as $u_1c_1+u_2c_2+u_3c_3$.
This agrees with the character indexing in that example. The nonzero
character supports and quotient genera are
\[
\begin{array}{c|c|c}
 \text{character word}&\text{active support}&\text{quotient genus}\\ \hline
 c_2&\{1,2,3,4,5,6\}&2\\
 c_1+c_2&\{1,2,3,6\}&1\\
 c_2+c_3&\{1,2,3,4\}&1\\
 c_1+c_2+c_3&\{1,2,3,5\}&1\\
 c_1&\{4,5\}&0\\
 c_3&\{5,6\}&0\\
 c_1+c_3&\{4,6\}&0
\end{array}
\]
The full cover has genus five, also directly from
$2g-2=8(-2+6\cdot\frac12)=8$.

The pair twists at $(1,6)$, $(1,4)$, and $(1,5)$ act nontrivially on,
respectively, the three displayed elliptic quotient factors and
projectively trivially on the other two. The nonseparating vanishing-cycle
argument of Section~\ref{sec:binary} and the graph normalizer therefore
make their three $\SL_2$ factors independent. The genus-two factor has
group $\Sp_4$ and is independent by the different rank. Hence, as
rational groups in the indicated character representation,
\[
 M=\Sp_4\times\SL_2^3
       =\Sp(H^1(JX,\Q))^G.
\]
At a very general point the inclusions
$M\subset\operatorname{Hg}(JX)\subset\Sp(H^1(JX,\Q))^G$
are therefore equalities. If $S_f$ denotes the smallest special
subvariety containing the family and $S(G)$ its deck-defined PEL
subvariety, on the component containing the family this proves
\[
                         S_f=S(G),\qquad\dim S_f=3+1+1+1=6.
\]
The branch family has dimension three and is not special.
The equality $S_f=S(G)$ is the issue left undecided in the cited
example; the support argument supplies it. It is also a short consequence
of older results: Bor\'owka--Ortega \cite[Remark 3.5, Proposition 3.9]{BO}
already describe the same three independently varying Legendre factors
for the non-isotropic Klein cover of the genus-two quotient. Their
decomposition, classical genus-two monodromy and Goursat give the same
answer. Thus resolving that precise formulation does not establish a new
monodromy phenomenon or earliest priority for this corollary.

\section{Dependencies and verification boundary}\label{sec:boundary}
The arguments supplied here are the mixed-embedding carry lemma,
monodromy support and labelled-pair reconstruction, oriented graph
construction, rational diagram realization, and Hodge lifting to balanced
sources. The following inputs have distinct roles:
\begin{itemize}
\item Menet--Nguyen, Definition~1.1(a), Theorems~2.1, 2.2, 2.5, 5.1,
and Corollary~2.7: cyclic dimensions, pair operators, and full factor
projections. All hypotheses used are matched in Sections~2--3.
\item The companion finite-quotient construction \cite{Fermat} and
Shioda's Theorem~2 \cite{Shioda}: determinant spaces are algebraic direct
summands of Fermat cohomology, and Hodge classes on their products are
algebraic. The common singular quotient is handled by a Chow-correspondence
identity between smooth varieties. This route needs no Aoki pairing,
admissible-cover smoothing, or cycle specialization.
\item Classical automorphisms of adjoint linear groups, semisimple Goursat,
Schur--Weyl duality and symplectic tensor invariants. The particular mixed
$\SL$ tensor statement is proved from Schur--Weyl in
Lemma~\ref{lem:fft}; no finite computation replaces it.
\item Generic monodromy and algebraic Hodge loci, with the polarization
complement argument for Hodge semisimplicity.
\item A'Campo and Picard--Lefschetz for the binary case.
\end{itemize}

The repository's exact finite arithmetic checks exercise the carry formula,
rank-two sign cube, higher-rank reconstruction, and several false shortcuts.
They are regression evidence, not proofs of universal quantifiers or of
cycle algebraicity. Its existing Lean results cover only designated
arithmetic and abstract interfaces; this entire geometric argument is not
formalized. The separate companion \cite{Extensions} supplies a new
geometric centralizer proof for the exact generic Hodge group; it does not
rehabilitate the erroneous legacy formulas. It also treats products across
distinct primes and explicit moving absolutely simple Jacobians whose first
exceptional classes and all-powers generators can be determined.
Its Section 7 supplies the broader odd-exponent criterion and unconditional
Aoki-exponent cases using primitive cyclic factors and controlled Fermat
blowups; its Section 6 gives the degree-seven generalized-Hodge corollary.
Independent specialist review and exact novelty assessment remain necessary
before treating these manuscripts as publication-ready.

\begin{thebibliography}{99}
\bibitem{MN} G. Menet and D.-M. Nguyen,
\emph{Representations of braid groups via cyclic covers of the sphere:
Zariski closure and arithmeticity}, arXiv:2310.10401v3.
\url{https://arxiv.org/html/2310.10401v3}.
\bibitem{Fusion} K. Kistner, \emph{Balanced mixed determinants of prime
cyclic covers: a revised fusion argument}, 4 September 2026.
Companion source \texttt{manuscripts/fusion\_revised.tex},
Theorem~1.1 and Sections~2--8.
\bibitem{Fermat} K. Kistner, \emph{Determinants of cyclic covers of the line:
a finite-quotient reduction to Fermat varieties}, 5 September 2026.
Companion source \nolinkurl{manuscripts/fermat_determinant_transfer.tex}.
\bibitem{Extensions} K. Kistner, \emph{Hodge groups, exceptional classes,
and odd-exponent abelian covers},
5 September 2026. Companion source
\nolinkurl{manuscripts/structural_extensions.tex}.
\bibitem{Shioda} T. Shioda, \emph{The Hodge conjecture and the Tate
conjecture for Fermat varieties}, Proc. Japan Acad. Ser. A Math. Sci.
\textbf{55} (1979), 111--114, Theorem~2, p.~112.
\url{https://doi.org/10.3792/pjaa.55.111}.
\bibitem{Rohde} J. C. Rohde, \emph{Cyclic coverings, Calabi--Yau manifolds
and Complex multiplication}, dissertation, Duisburg--Essen, 2007,
Construction 3.2.1, Section 4.3, Theorem 5.3.7.
\url{https://webdoc.sub.gwdg.de/ebook/dissts/Duisburg/Rohde2007.pdf}.
\bibitem{ST} I. Spelta and C. Tamborini, \emph{Decomposable abelian
$G$-curves and special subvarieties}, Bulletin des Sciences Math\'ematiques
201 (2025), 103616, Theorem 1.7 and Section 3.
\url{https://doi.org/10.1016/j.bulsci.2025.103616}.
\bibitem{FMZ} G. Forni, C. Matheus and A. Zorich,
\emph{Square-tiled cyclic covers}, Journal of Modern Dynamics
\textbf{5} (2011), 285--318, Section 2.6, Lemma 16.
\url{https://doi.org/10.3934/jmd.2011.5.285}.
\bibitem{BO} P. Bor\'owka and A. Ortega, \emph{Klein coverings of genus 2
curves}, Transactions of the American Mathematical Society
\textbf{373} (2020), 1885--1907. Author manuscript, Remark 3.5 and
Proposition 3.9, \url{https://arxiv.org/pdf/1904.05962v2}.
\bibitem{Andre} Y. Andr\'e, \emph{Mumford--Tate groups of mixed Hodge
structures and the theorem of the fixed part}, Compositio Mathematica
\textbf{82} (1992), 1--24, Section~4, Lemma~4.
\url{https://www.numdam.org/article/CM_1992__82_1_1_0.pdf}.
\bibitem{CDK} E. Cattani, P. Deligne and A. Kaplan,
\emph{On the locus of Hodge classes}, Journal of the American Mathematical
Society \textbf{8} (1995), 483--506, Theorem~1.1 and Corollary~1.2.
\url{https://arxiv.org/abs/alg-geom/9402009}.
\bibitem{ACampo} N. A'Campo, \emph{Tresses, monodromie et le groupe
symplectique}, Commentarii Mathematici Helvetici \textbf{54} (1979),
318--327, Theorem~1.
\url{https://doi.org/10.1007/BF02566275}.
\bibitem{Etingof} P. Etingof, O. Golberg, S. Hensel, T. Liu, A. Schwendner,
D. Vaintrob and E. Yudovina, \emph{Introduction to representation theory},
author manuscript, 10 January 2011, Sections~4.18--4.19, p.~65.
\url{https://math.mit.edu/~etingof/replect.pdf}.
\bibitem{Goodman} R. Goodman, \emph{Multiplicity-free spaces and
Schur--Weyl--Howe duality}, revised 3 July 2003, Section~7.3,
Theorem~7.4 and Corollary~7.5, pp.~26--28.
\url{https://sites.math.rutgers.edu/~goodman/pub/nuslect.pdf}.
\end{thebibliography}
\end{document}
